%============================================================
% HALAMAN PENGESAHAN
% Watermark background otomatis untuk 1 atau lebih halaman.
%============================================================
\cleardoublepage
\pagestyle{empty} % Memastikan seluruh halaman pengesahan tanpa nomor/header
\addtocontents{toc}{\protect\cftpagenumbersoff{chapter}}
\addcontentsline{toc}{chapter}{HALAMAN {PENGESAHAN}}
\addtocontents{toc}{\protect\cftpagenumberson{chapter}}

\ifpengesahanpdf
% Sisipkan lembar pengesahan PDF yang sudah ditandatangani
% (aktifkan \pengesahanpdftrue di main_buku.tex / config/toggles.tex)
\includepdf[pages=-]{\FilePengesahanPDF}
\else

% 1. Nyalakan watermark secara otomatis mulai dari halaman ini
\AddToShipoutPictureBG{%
  \tikz[remember picture, overlay]
  \node[inner sep=0pt] at (current page.center){%
    \includegraphics[width=\paperwidth, height=\paperheight]{assets/pens/background_pens_repeated.png}%
  };
}

\begin{center}
  {\FontPengesahanLabel\bfseries HALAMAN PENGESAHAN\par}
  \vspace{0.8cm}

  \begin{minipage}{1\textwidth}
    \centering
    \FontPengesahanJudul\bfseries\MakeUppercase{\JudulPA}
  \end{minipage}
  \par

  \vspace{0.6cm}

  {\FontPengesahanIsi Oleh:\par}
  \vspace{0.1cm}

  {\FontPengesahanIsi\bfseries
    \NamaMahasiswa\par
  NRP. \NRP\par}

  \vspace{0.6cm}

  \begin{spacing}{1}
    {\FontPengesahanIsi\bfseries
      Proyek Akhir ini digunakan sebagai salah satu syarat untuk\\
      memperoleh Gelar \GelarLulus\\
      di Politeknik Elektronika Negeri Surabaya\\
    \TahunLulus\par}
  \end{spacing}
\end{center}

\vspace{0.2cm}
\begin{center}
  {\FontPengesahanIsi\bfseries Disetujui oleh:}
\end{center}
\vspace{0.2cm}

% Daftar pembimbing dan penguji beserta kolom tanda tangan
\noindent
{\FontPengesahanIsi
  \begin{tabular}{@{}l c l c@{}}
    Pembimbing 1 & : & \NamaPembimbingSatu   & (\hspace{2.5cm}) \\
    &   & NIP. \NIPPembimbingSatu & \\[3pt]
    Pembimbing 2 & : & \NamaPembimbingDua   & (\hspace{2.5cm}) \\
    &   & NIP. \NIPPembimbingDua & \\[3pt]

    % Baris kondisional Pembimbing 3 (Akademik)
    \ifcsvoid{NamaPembimbingTiga}{}{%
      Pembimbing 3 & : & \NamaPembimbingTiga   & (\hspace{2.5cm}) \\
      &   & NIP. \NIPPembimbingTiga & \\[3pt]
    }

    % Baris kondisional Pembimbing Industri
    \ifcsvoid{NamaPembimbingIndustri}{}{%
      Pembimbing Industri & : & \NamaPembimbingIndustri & (\hspace{2.5cm}) \\
      &   &                        & \\[3pt]
    }

    Penguji 1    & : & \NamaPengujiSatu     & (\hspace{2.5cm}) \\
    &   & NIP. \NIPPengujiSatu   & \\[3pt]
    Penguji 2    & : & \NamaPengujiDua      & (\hspace{2.5cm}) \\
    &   & NIP. \NIPPengujiDua    & \\
  \end{tabular}
}

\vspace{0.3cm}

% 2. Blok Mengetahui dibungkus dalam minipage
%    Jika ruang tidak cukup, SELURUH kotak ini akan pindah ke halaman 2.
\noindent
\begin{minipage}{\textwidth}
  \begin{center}
    {\FontPengesahanIsi\bfseries Mengetahui,}\par
    {\FontPengesahanIsi
      Ketua Program Studi Sarjana Terapan \NamaProdi\par
    Politeknik Elektronika Negeri Surabaya}\par
    \vspace{1.5cm}
    {\FontPengesahanIsi
      \underline{\NamaKetuaProdi}\par
    NIP. \NIPKetuaProdi\par}
  \end{center}
\end{minipage}

% 3. Paksa render halaman terakhir, lalu MATIKAN watermark agar tidak bocor
\clearpage
\ClearShipoutPictureBG
\fi